\documentclass[11pt,a4paper]{article}
\usepackage[T1]{fontenc}
\usepackage[utf8]{inputenc}
\usepackage{amsmath,amssymb,xcolor}
\usepackage[margin=25mm]{geometry}
\usepackage[colorlinks=true,urlcolor=blue]{hyperref}
\definecolor{nexorablue}{HTML}{2459B5}
\setlength{\parindent}{0pt}
\setlength{\parskip}{7pt}
\setlength{\emergencystretch}{2em}
\renewcommand{\familydefault}{\sfdefault}
\begin{document}

{\small\bfseries\color{nexorablue} NEXORA CPGE}\par
{\small MPSI / NEX-CIR-03 / Corrigé}\par
{\LARGE\bfseries Une sortie avec et sans charge}\par
\vspace{0.25em}\hrule\vspace{0.6em}

\section*{Sortie sans charge}
Les valeurs en volts et kilo-ohms donnent $(12-V_B)/2=V_B/4$, d’où $V_B=8\,\mathrm V$. Le courant vaut $2\,\mathrm{mA}$.

\section*{Sortie en charge}
Avec la nouvelle branche, $(12-V_B)/2=V_B/4+V_B/4$. On obtient $V_B=6\,\mathrm V$. Le courant entrant en $B$ vaut $3\,\mathrm{mA}$. Il se partage en deux courants de $1{,}5\,\mathrm{mA}$, dirigés vers la masse.


\end{document}
